\documentclass[10pt,a4paper]{article}
\usepackage[T1]{fontenc}\usepackage[utf8]{inputenc}\usepackage[english,italian]{babel}
\usepackage{lmodern,microtype}\usepackage[a4paper,margin=1.85cm,headheight=15pt]{geometry}
\usepackage{enumitem,tabularx,array,longtable,booktabs,pdflscape,xcolor,hyperref,xurl,fancyhdr,titlesec,framed,listings}
\definecolor{ddtadark}{RGB}{28,38,52}\definecolor{ddtablue}{RGB}{42,88,135}\definecolor{ddtagray}{RGB}{244,246,248}\definecolor{ddtagreen}{RGB}{48,111,78}\definecolor{ddtaorange}{RGB}{148,91,25}
\hypersetup{colorlinks=true,linkcolor=ddtablue,urlcolor=ddtablue}
\pagestyle{fancy}\fancyhf{}\lhead{DDTA - BA0-R}\rhead{Research history R3}\cfoot{\thepage}
\titleformat{\section}{\large\bfseries\color{ddtadark}}{\thesection}{.6em}{}\titleformat{\subsection}{\normalsize\bfseries\color{ddtadark}}{\thesubsection}{.6em}{}
\setlist[itemize]{leftmargin=1.45em,itemsep=.16em,topsep=.2em}\setlist[enumerate]{leftmargin=1.7em,itemsep=.16em,topsep=.2em}
\setlength{\parindent}{0pt}\setlength{\parskip}{.32em}\setlength{\emergencystretch}{2em}
\lstset{basicstyle=\ttfamily\small,frame=single,breaklines=true,backgroundcolor=\color{ddtagray}}
\newenvironment{challenge}[1]{\def\FrameCommand{\fcolorbox{ddtaorange}{ddtagray}}\MakeFramed{\advance\hsize-2\width\FrameRestore}\noindent\textbf{#1}\par\smallskip}{\endMakeFramed}
\newenvironment{gate}[1]{\def\FrameCommand{\fcolorbox{ddtagreen}{ddtagray}}\MakeFramed{\advance\hsize-2\width\FrameRestore}\noindent\textbf{#1}\par\smallskip}{\endMakeFramed}
\hypersetup{pdftitle={BA0-R systems-modeling prior-art synthesis Revision 3}}
\begin{document}
\begin{center}{\LARGE\bfseries BA0-R systems-modeling prior-art synthesis}\par{\large Revision 3 - research history}\par\smallskip{\small Baseline: \texttt{89e5486e02d07ff5b97082f73a2e22ac1b4319ae}}\end{center}
\begin{challenge}{Status}Includes SRC-0048, the full-text ISO/IEC/IEEE 42010:2011 regression. Semantic closure remains OPEN; no BAE metaclass is promoted.\end{challenge}
\section{Corpus and access boundary}
Revision 3 preserves the R2 corpus and adds SRC-0048, a user-supplied full-text reading of ISO/IEC/IEEE 42010:2011. SRC-0043 remains the separate ACCESS-LIMITED 2022 record; no equivalence between editions is assumed.
\section{Material correction from ISO 42010:2011}
The earlier direction \texttt{canonical model -> projection -> view} was too strong as a statement of general prior art. ISO 42010 describes both synthetic construction of views integrated through correspondences and projective extraction of views from an underlying repository.
\begin{lstlisting}
one canonical repository is universally mandatory   REJECTED as prior-art claim
projective DDTA Base Analysis                       CANDIDATE design choice
\end{lstlisting}
DDTA may still choose one accepted Base Analysis to simplify consistency, provenance, refresh and re-analysis, but BA0/BA4 must justify that choice.
\section{View terminology}
\begin{tabularx}{\linewidth}{>{\bfseries}p{3.5cm}X}
Viewpoint & reusable conventions framing concerns and specifying model kinds plus construction/interpretation/analysis methods.\\
View & application of one viewpoint to one system-of-interest; its architecture models collectively address the framed concerns across the whole system.\\
Architecture Model & model governed by a model kind; may selectively present part of the system and may be shared by multiple views.\\
\end{tabularx}
A partial diagram or filtered graph should not automatically be called an architecture view if DDTA adopts ISO-aligned terminology.
\section{ISO 42010 does not solve BA1}
The standard explicitly does not define what a system is. Its ontology governs architecture-description work products, not the internal system-model ontology. Annex B nevertheless provides a useful metamodel lens: entities, attributes, relationships and constraints.
\section{Cross-source responsibility status}
\begin{landscape}\scriptsize
\begin{longtable}{p{4cm}p{10cm}p{8cm}}
\toprule Responsibility & Evidence after NASA/SysML/KerML/AADL/OPM/ISO & DDTA status\\\midrule\endhead
Structure / referent & strong across traditions & semantic responsibility STRONG; exact taxonomy OPEN\\
Behavior / transformation & function/action/behavior/Process evidence & STRONGLY EVIDENCED\\
Information/resource & recurring but represented differently & semantic need strong; universal InformationObject root OPEN/WEAKENED\\
Relation/connection & connector/link/correspondence evidence & typed relation semantics STRONGLY EVIDENCED\\
Interface/endpoint & recurrent, often viewpoint/model-kind specific & first-class identity OPEN\\
Flow/transfer & recurrent; can be relation semantics & root OPEN/WEAKENED\\
State/mode & recurrent; can be owned/classified & root OPEN/WEAKENED\\
Boundary/environment & distinction strong, root not forced & scope distinction STRONG; Boundary root OPEN/WEAKENED\\
Property/constraint & explicit across metamodels & likely property/constraint mechanism\\
View/viewpoint/model kind & direct ISO full-text support & distinctions STRONGLY EVIDENCED\\
Cross-model relations & ISO correspondences plus other approaches & STRONGLY EVIDENCED prior art\\
Analysis & SysML/AADL/NASA/ISO & generic model-based analysis is NOT DDTA novelty\\
\bottomrule\end{longtable}\end{landscape}
\section{Hypothesis changes}
Actor root and Asset root remain REJECTED as generic roots. Interaction remains REJECTED as a sufficient universal relation. Channel, Store, Contract, Boundary, State, Flow and InformationObject remain OPEN/WEAKENED as mandatory roots. ISO additionally rejects the claim that a single repository is required by architecture-description prior art.
\section{Correspondence and provenance pressure}
ISO correspondences can be explicit n-ary relations governed by rules and used for consistency, traceability, composition, refinement, dependency and model transformation, including links between architecture descriptions and requirements. DDTA historical analysis-to-governed-change provenance remains a distinct unresolved problem, but BA3/BA4 must compare against correspondence prior art.
\section{Decision and rationale regression}
Architecture decisions and rationale, including affected elements, alternatives, consequences, timestamps and citations, are established prior art. DDTA Decision remains distinct because its closed responsibility is normative ownership/decomposition under one MacroRequirement. No Chapter 4 reopen trigger is produced.
\section{Novelty challenge}
Established prior art now includes modeling metamodels, multiple views/viewpoints, model-based analysis, cross-model consistency/traceability and decision rationale. The candidate DDTA contribution must remain the governed composition boundary:
\begin{lstlisting}
governed documentation
 -> reviewed methodology-neutral analyzable representation
 -> replaceable method-owned security analysis
 -> reviewed normalized findings
 -> accepted governed document / requirement change
 -> preserved analysis-to-change provenance
 -> change-aware re-analysis
\end{lstlisting}
The reviewed systems-modeling corpus does not establish that entire loop. This is a bounded comparative inference, not proof of novelty.
\section{Chapter regression}
\begin{lstlisting}
Chapter 3 reopen trigger  NOT FOUND
Chapter 4 reopen trigger  NOT FOUND
\end{lstlisting}
\section{BA0 working hypothesis - OPEN}
\begin{gate}{Working hypothesis}
Base Analysis is an accepted, methodology-neutral analyzable representation of governed project/system knowledge, grounded in governed documentation and explicit reviewed analytical additions. It must preserve analysis-relevant semantic distinctions without adopting one external modeling-language ontology. DDTA may choose a canonical projective representation to support repeatable views, consistency, provenance and re-analysis, but that choice must be justified against synthetic multi-view alternatives rather than assumed as a universal modeling principle.
\end{gate}
\section{Next decision}
After the history package is committed, decide whether to merge the candidate sources into the canonical literature registry and close BA0-R. No BAE metaclass is accepted here.
\end{document}
